\chapter{Damping}
\label{ch:damping}

\begin{idea}
	Damping needs a resistance of order $R_0$ in the resonant loop that carries
	no DC. Putting $R_0$ in series with the line would burn power and drop
	voltage. A damping \emph{leg} avoids that:
	\begin{itemize}
		\item An $R_d$--$L_d$ leg across $L$ sees no DC, because $L$ shorts it.
		\item An $R_d$--$C_d$ leg across $C$ sees no DC, because $C_d$ blocks it.
	\end{itemize}
\end{idea}

The design question for either leg is how big to make it. Call its size
$n$: $n = L_d/L$ or $n = C_d/C$. For each $n$ there is one $R_d$ that
minimizes the worst-case $\Zo$ peak. That optimum peak is written
\[
	\norm{\Zo}_{mm} = \min_{R_d}\left(\max_f \abs{\Zo}\right).
\]

All closed forms here are from Erickson \& Maksimovi\'c, Ch.~17
\cite{erickson}, and were checked against a brute-force optimum to 4 digits.
$R_0 = \sqrt{L/C}$ is always the \emph{undamped} section's value.

% -------------------------------------------------------------------

\section{\texorpdfstring{$R_d$--$L_d$}{Rd-Ld} leg across the inductor}
\label{sec:rd-ld}

\subsection{What the loop sees}

\begin{derivation}
	The inductor with its leg is $sL \parallel (R_d + sL_d)$:
	\[
		Z = \frac{sL(R_d + sL_d)}{R_d + s(L + L_d)}.
	\]
	At $s = j\omega$, multiply through by the conjugate of the denominator
	and keep the real part:
	\[
		\operatorname{Re}\left[(j\omega L R_d - \omega^2 L L_d)(R_d - j\omega(L + L_d))\right]
		= -\omega^2 L L_d R_d + \omega^2 L R_d (L + L_d) = \omega^2 L^2 R_d.
	\]
	Dividing by $\abs{R_d + j\omega(L+L_d)}^2$:
	\begin{equation}\label{eq:reff}
		R_{eff}(\omega) = \frac{R_d(\omega L)^2}{R_d^2 + \omega^2(L + L_d)^2}.
	\end{equation}
\end{derivation}

So the leg adds a series resistance that is zero at DC and turns on above the corner
\[
	\omega_c = \frac{R_d}{L + L_d}, \qquad
	R_{eff} \to \frac{R_d L^2}{(L + L_d)^2} = \frac{R_d}{(1+n)^2}\ \text{above } \omega_c.
\]

\subsection{The cost: lost high-frequency inductance}

Well above $\omega_c$, $sL_d \gg R_d$ and the leg looks like a plain inductor in
parallel with $L$:
\begin{equation}\label{eq:lhf}
	L_{HF} = L \parallel L_d = \frac{nL}{1+n}, \qquad
	\text{penalty}_{dB} = 20\log_{10}\frac{1+n}{n}.
\end{equation}

To keep a target $L_{HF}$, you need $L = L_{HF}(1+n)/n$. Small $n$ damps hard
and costs a lot of attenuation. $n = 1$ costs 6~dB.

\subsection{The optimum, by the fixed-point method}

\begin{idea}
	Look at the two extremes of $R_d$:
	\begin{itemize}
		\item $R_d \to 0$: a lossless tank of $(L \parallel L_d)$ with $C$.
		\item $R_d \to \infty$: a lossless tank of $L$ with $C$.
	\end{itemize}
	Those two $\abs{\Zo}$ curves cross at one frequency, and \emph{every}
	$R_d$ curve passes through that crossing. No $R_d$ can get the peak below
	the crossing height. The optimum $R_d$ puts the peak exactly there, with zero slope.
\end{idea}

\begin{derivation}
	Write $L' = L \parallel L_d = nL/(1+n)$. A lossless tank $sL \parallel 1/(sC)$ has
	\[
		\Zo = \frac{j\omega L}{1 - \omega^2 LC}.
	\]
	At the crossing, the $R_d = 0$ tank is below its resonance and the
	$R_d = \infty$ tank is above its own, so the two have equal magnitude and
	opposite sign:
	\[
		\frac{\omega L}{1 - \omega^2 LC} = -\frac{\omega L'}{1 - \omega^2 L'C}
		\quad\Rightarrow\quad
		L + L' = 2\omega^2 L L' C
		\quad\Rightarrow\quad
		\omega^2 = \frac{L + L'}{2LL'C}.
	\]
	Substituting $L' = nL/(1+n)$ gives $L + L' = L(1+2n)/(1+n)$ and
	$LL' = nL^2/(1+n)$, so
	\[
		\omega^2 = \frac{1 + 2n}{2n\,LC}.
	\]
	The height there, using $1 - \omega^2 LC = 1 - \frac{1+2n}{2n} = -\frac{1}{2n}$:
	\[
		\abs{\Zo} = \frac{\omega L}{1/(2n)} = 2n\,\omega L
		= 2n R_0 \sqrt{\frac{1+2n}{2n}}
		= R_0\sqrt{2n(1+2n)}.
	\]
\end{derivation}

Setting $d\abs{\Zo}^2/d\omega = 0$ at that crossing and solving for $R_d$
gives the optimum resistor. The algebra is long; the results are
\begin{equation}\label{eq:rd-ld-opt}
	\norm{\Zo}_{mm} = R_0\sqrt{2n(1+2n)}, \qquad
	R_{d,opt} = R_0\sqrt{\frac{n(3+4n)(1+2n)}{2(1+4n)}}.
\end{equation}

\subsection{Choosing \texorpdfstring{$n$}{n}}

Pick the \emph{largest} $n$ that meets the $\Zo$ target $Z_t$, since
over-damping costs attenuation you do not get back.

\begin{derivation}
	Set the optimum peak equal to the target and let $k = (Z_t/R_0)^2$:
	\[
		2n(1+2n) = k \quad\Rightarrow\quad 4n^2 + 2n - k = 0
		\quad\Rightarrow\quad
		n_{max} = \frac{-1 + \sqrt{1 + 4k}}{4}.
	\]
\end{derivation}

Parasitic DCR and ESR add to $R_{eff}$, so the exact peak sits below
\zcref{eq:rd-ld-opt}. The formula is conservative, and the gap is largest
at large $n$, where $R_d$ itself contributes less.

\begin{figure}[ht]
	\centering
	\begin{tikzpicture}
		\begin{loglogaxis}[
				width=0.9\textwidth, height=6.5cm,
				xmin=2e3, xmax=4e4, ymin=0.5, ymax=120,
				xtick={2e3,5e3,1e4,2e4,4e4}, xticklabels={2k,5k,10k,20k,40k},
				legend style={font=\small},
				xlabel={Frequency (Hz)}, ylabel={$\abs{\Zo}$ ($\Omega$)},
				legend pos=north east, legend cell align=left,
				grid=both, grid style={black!8},
			]
			\addplot[black, thick] table[x=f, y=undamped] {figures/single-stage-zo.dat};
			\addplot[blue, thick] table[x=f, y=n05] {figures/single-stage-zo.dat};
			\addplot[green!50!black, thick] table[x=f, y=n1] {figures/single-stage-zo.dat};
			\addplot[orange, thick] table[x=f, y=nmax] {figures/single-stage-zo.dat};
			\addplot[red, dashed, thick, domain=2e3:4e4] {26.9};
			\legend{undamped, {$n=0.5$}, {$n=1$}, {$n=4$}, {$\abs{\Zin}/3.16$}}
		\end{loglogaxis}
	\end{tikzpicture}
	\caption{The \zcref{ex:first-try} stage, undamped and with an optimally sized
		$R_d$--$L_d$ leg. Larger $n$ leaves a higher peak but costs less attenuation.}
	\label{fig:single-stage-zo}
\end{figure}

% -------------------------------------------------------------------

\section{\texorpdfstring{$R_d$--$C_d$}{Rd-Cd} leg across the capacitor}
\label{sec:rd-cd}

Same idea, with the leg $R_d + 1/(sC_d)$ across $C$ and $n = C_d/C$.
This leg costs no high-frequency attenuation (above its corner $C_d$ just adds
to $C$), but $C_d$ is usually a large capacitor.

\begin{derivation}
	The extremes: $R_d \to 0$ gives $L$ with $C(1+n)$; $R_d \to \infty$ gives
	$L$ with $C$. Equal magnitude, opposite sign:
	\[
		1 - \omega^2 LC(1+n) = -(1 - \omega^2 LC)
		\quad\Rightarrow\quad
		\omega^2 = \frac{2}{(2+n)LC}.
	\]
	Then $1 - \omega^2 LC = n/(2+n)$, and the height is
	\[
		\abs{\Zo} = \frac{\omega L}{n/(2+n)} = R_0\sqrt{\frac{2}{2+n}}\cdot\frac{2+n}{n}
		= R_0\frac{\sqrt{2(2+n)}}{n}.
	\]
\end{derivation}

\begin{equation}\label{eq:rd-cd-opt}
	\norm{\Zo}_{mm} = R_0\frac{\sqrt{2(2+n)}}{n}, \qquad
	R_{d,opt} = R_0\sqrt{\frac{(2+n)(4+3n)}{2n^2(4+n)}}.
\end{equation}

Here a bigger leg damps harder, so you want the \emph{smallest} $n$ that
meets the target. With $k = (Z_t/R_0)^2$, $kn^2 - 2n - 4 = 0$ gives
\[
	n_{min} = \frac{1 + \sqrt{1 + 4k}}{k}.
\]

\begin{examplebox}[label=ex:damp-first-try]{damping the first-try LC}
	Add an $R_d$--$L_d$ leg to the \zcref{ex:first-try} stage
	($R_0 = \qty{3.16}{\ohm}$). The target from \zcref{ex:zin} is
	$Z_t = \qty{26.9}{\ohm}$, so
	\[
		k = \left(\frac{26.9}{3.16}\right)^2 = 72.4, \qquad
		n_{max} = \frac{-1 + \sqrt{1 + 4(72.4)}}{4} = 4.0.
	\]
	\zcref[S]{eq:rd-ld-opt} and \zcref{eq:lhf} for three choices of $n$, with the
	exact results from a sweep that includes the \qty{0.13}{\ohm} of parasitics:

	\begin{center}
		\small
		\begin{tabular}{@{}lllllll@{}}
			\toprule
			$n$   & $R_{d,opt}$      & Formula peak     & Exact peak       & Margin  & Penalty & $\abs{H}$ at \qty{400}{\kHz} \\
			\midrule
			0.5   & \qty{2.89}{\ohm} & \qty{4.47}{\ohm} & \qty{4.32}{\ohm} & 25.9~dB & 9.5~dB  & $-59.6$~dB                   \\
			1     & \qty{4.58}{\ohm} & \qty{7.75}{\ohm} & \qty{7.22}{\ohm} & 21.4~dB & 6.0~dB  & $-63.2$~dB                   \\
			4.0   & \qty{14.2}{\ohm} & \qty{26.9}{\ohm} & \qty{20.4}{\ohm} & 12.4~dB & 1.9~dB  & $-67.2$~dB                   \\
			\bottomrule
		\end{tabular}
	\end{center}

	Every row is now stable, but look at the cost. $n = 0.5$ buys 26~dB of
	stability margin, far more than needed, and gives up 9.5~dB of attenuation
	for it. $n = 4$ meets the target while giving up only 1.9~dB.
	The exact peaks sit below the formula because DCR and ESR damp on top of
	$R_d$, most noticeably at large $n$ where $R_d$ adds the least.

	None of these meets the 80.4~dB requirement, since the undamped stage only
	managed 69.2~dB to begin with.
\end{examplebox}

\begin{examplebox}[label=ex:single-damped]{a single damped stage that meets the requirement}
	\zcref[S]{ex:requirement} needs $L = \qty{244}{\micro\henry}$ with
	\qty{6.8}{\micro\farad} for 80.4~dB. With a damping leg that is the
	inductance that must be left at high frequency, $L_{HF}$, so
	$L = L_{HF}(1+n)/n$.

	With $n = 1$:
	\[
		L = L_d = 2(244) = \qty{489}{\micro\henry}, \qquad
		R_0 = \sqrt{\frac{489}{6.8}}\,\Omega = \qty{8.48}{\ohm}, \qquad
		\norm{\Zo}_{mm} = 8.48\sqrt{6} = \qty{20.8}{\ohm}.
	\]
	That passes the \qty{26.9}{\ohm} target. With $n = 0.5$ it takes
	\qty{733}{\micro\henry} plus \qty{367}{\micro\henry}. Going the other way,
	$n = 1.5$ just passes (\qty{26.8}{\ohm}) and $n = 2$ fails (\qty{32.8}{\ohm}).

	Either way the stage needs about \qty{1}{\milli\henry} of total inductance.
	The two-stage design in \zcref{ch:two-stage} meets the same requirement with
	about \qty{83}{\micro\henry} and the same total capacitance.
\end{examplebox}
